\section*{v1.3.0 Expansion Core}
\addcontentsline{toc}{section}{v1.3.0 Expansion Core}

\begin{remark}[Normalization note for v1.3.0]
\label{rem:normalization-note-v130}
This block is retained as a historical expansion precursor. Its proof-bearing content is now to be read primarily through Phase~I, Phase~II-D, and the normalized grouped architecture. In particular, the intermediate master-equation ladder and the micro-to-macro chain no longer function as independent foundations beside the later unfolding phases, but as early organizational anticipations of them.
\end{remark}

Version v1.3.0 is the first release that does more than freeze the next
work path. It adds a genuine expansion layer on top of the v1.1.0
consolidation core. The operative aim is threefold: first, the structural
master equation is no longer presented only as a compressed global summary, but
also as the endpoint of an explicit intermediate chain; second, the passage
from local triolic organization to macroscopic effective geometry is stated in
staged form; third, the technical apparatus is widened so that later proof
layers can be attached without rewriting the monolith.

\subsection*{1. Deepened master-equation programme}
The master equation in the later main section remains the global compression of
one common closure-stable dynamics. What v1.3.0 adds is a more explicit
three-level reading of that equation.
\begin{enumerate}[leftmargin=2em,label=(\roman*)]
  \item \textbf{Structural level.} One starts from triolic organization,
  orthogonality, frame shift, compensator action, and closure. At this level,
  the equation is still read as a balance law on admissible sectors.
  \item \textbf{Mesoscopic level.} The same balance law is rewritten as an
  effective transport-and-selection equation for coupled sectors and probes.
  Here one already sees how admissible interaction channels are filtered by HC
  and by the closure residue.
  \item \textbf{Macroscopic level.} After suppression of unstable local
  residues and after passage to a closure-stable observer description, the same
  equation becomes readable as effective matter dynamics plus background
  geometry.
\end{enumerate}
In this sense, v1.3.0 does not replace the later master equation
\eqref{eq:toe_master_structural}; it inserts the missing ladder by which that
equation is reached.

A convenient intermediate schematic form is
\[
\mathcal E_{\mathrm{micro}}[\Psi]
+\mathcal E_{\mathrm{coupl}}[\Psi,\Pi]
+\mathcal E_{\mathrm{HC}}[\Psi]
+\mathcal E_{\mathrm{cl}}[\Psi]
+\mathcal E_{\mathrm{obs}}[\Psi]
=0,
\]
where the five blocks stand respectively for local triolic structure, sectorial
coupling/probe interaction, Heisenberg-compensator balancing, closure
selection, and observer-frame projection. The later master equation is to be
read as the compressed global representative of this staged balance.

\begin{lemma}[Intermediate balance principle]
\label{lem:intermediate-balance}
Every admissible global solution must pass through an intermediate regime in
which local triolic transport, compensator balancing, and closure selection are
all simultaneously active. Consequently, a direct jump from isolated local
structure to macroscopic field form is structurally forbidden.
\end{lemma}

\begin{proof}
Local triolic data alone do not yet determine which coupled configurations are
admissible; they only provide possible sectorial directions. Conversely, a
macroscopic field description is meaningful only after compensator balancing
has removed representation drift and closure has excluded unstable residues.
Therefore at least one intermediate regime must exist in which coupling, HC,
and closure are jointly active. This regime is the structural source of the
mesoscopic reading introduced in v1.3.0.
\end{proof}

\subsection*{2. Micro-to-macro closure chain}
The new hardening target of v1.3.0 is the explicit micro-to-macro chain
\[
\begin{aligned}
&\text{local triolic unit} \longrightarrow \text{coupled closure cell} \longrightarrow \text{probe-filtered admissible channel}\\
&\longrightarrow \text{closure-stable sector regime} \longrightarrow \text{effective global geometry and field dynamics}.
\end{aligned}
\]
The purpose of this chain is not rhetorical. It separates five logically
distinct burdens that earlier drafts could too easily blend together:
local decomposition, coupling, admissibility filtering, stable sector assembly,
and macroscopic readout.

\paragraph{Local triolic unit.}
At the smallest relevant level, the system carries the standing triolic
organization together with the orthogonal third direction. This level fixes the
minimal sector inventory but does not yet decide which transitions survive.

\paragraph{Coupled closure cell.}
Once two or more local units interact, the question becomes whether their
combination admits a common compensator-compatible representative. A coupled
cell is called closure-ready if its local mismatches can be balanced without
breaking the orthogonality and frame-shift discipline.

\paragraph{Probe-filtered admissible channel.}
The probe construction becomes technically important at precisely this stage.
Two directed probes together with the orthogonal third group define a restricted
search geometry. Kasisky-, hash-, rainbow-, Markov-, and Newton-type tools are
not external decoration here; they are the operational means by which the space
of candidate channels is compressed to closure-compatible paths.

\paragraph{Closure-stable sector regime.}
Only after probe filtering and compensator balancing does one obtain a stable
sector regime. This is the first level at which effective matter, light/collapse,
and dual sectors can be read as persistent rather than transient descriptions.

\paragraph{Effective global geometry and field dynamics.}
Macroscopic geometry is then not glued on from outside. It is the readout of a
large-scale assembly of closure-stable sector regimes viewed in the
phase-shifted observer frame. The gravitational, QFT-type, and Dirac-type
limits are therefore downstream realizations of one common filtered closure
architecture.

\begin{proposition}[Macroscopic readout requires closure-stable assembly]
\label{prop:macro-requires-closure-assembly}
A macroscopic field or geometric description is admissible only if it is the
observer-frame image of an assembly of closure-stable sector regimes. In
particular, large-scale geometry cannot be read directly from uncoupled local
triolic data.
\end{proposition}

\begin{proof}
Uncoupled local data do not yet fix admissible transport and do not suppress
representation drift. The observer frame, however, reads only stabilized
configurations. Since stabilization in the present framework is exactly the
combined action of HC and closure across coupled cells, the macroscopic image
must descend from an assembly that has already passed the closure filter.
\end{proof}

\subsection*{3. Extended lemma and tool apparatus}
Version v1.3.0 also widens the technical apparatus that the next proof
layer can rely on. The point is not to over-formalize prematurely, but to mark
where repeated arguments will later be detached into reusable blocks.

\paragraph{Core lemmas now fixed by role.}
\begin{enumerate}[leftmargin=2em]
  \item \textbf{Orthogonality lemma:} the third triolic direction remains the
  non-negotiable constraint that prevents degeneracy of the sector picture.
  \item \textbf{Frame-shift lemma:} observational readout differs from the
  underlying structural frame by a controlled phase displacement.
  \item \textbf{Compensator lemma:} HC is the minimal identification mechanism
  needed to prevent sectorial drift.
  \item \textbf{Closure lemma:} admissibility is equivalent to vanishing of the
  relevant closure residue in the chosen regime.
  \item \textbf{Probe admissibility lemma:} the probe sector selects channels
  that are compatible with both closure and orthogonality.
  \item \textbf{Read-off lemma:} physical sector names arise only after the
  preceding structural filters have been passed.
\end{enumerate}

\paragraph{Normalized tool classes.}
The tool catalogue of v1.1.0 is retained, but v1.3.0 upgrades its use:
closure-tools act as admissibility tests, frame-shift tools as translation
operators between structure and observer descriptions, projection-tools as
controlled reduction maps, probe-tools as search-space compressors,
equivalence-tools as representative-identification devices, and read-off tools
as disciplined output maps. This normalization matters because later technical
proofs can now cite a \emph{tool class} instead of repeating the surrounding
motivation every time.

\begin{corollary}[Tool usage discipline]
\label{cor:tool-usage-discipline}
No computational shortcut is admissible in the proof architecture unless its
role can be classified as one of the normalized tool classes and unless its use
preserves the structural ordering
\[
\text{closure before projection before interpretation}.
\]
\end{corollary}

\begin{proof}
A shortcut that bypasses closure may produce effective formulas without fixing
admissibility; a shortcut that bypasses projection discipline may confuse
structural content with observer naming. The normalized tool classes are
introduced exactly to block such category errors. Hence admissible shortcuts are
those that preserve the structural order stated above.
\end{proof}

\subsection*{4. Version status of v1.3.0}
With these additions, v1.3.0 should be read as the first expansion-stage
companion release. Relative to v1.1.0, the new version contributes three hard
improvements: the master-equation programme is given an explicit intermediate
ladder, the micro-to-macro chain is stated as a formal burden rather than a
mere promise, and the lemma/tool apparatus is widened enough that subsequent
versions can deepen proofs without reopening the global organization of the
monolith. The next versions may therefore focus on tightening derivations,
\[
\begin{aligned}
&\text{ground condition}
\Longrightarrow
\text{triolic admissibility}
\Longrightarrow
\text{HC balancing}
\\[0.2em]
&\Longrightarrow
\text{closure filtering}
\Longrightarrow
\text{observer-frame readout}.
\end{aligned}
\]
Version v1.4.0 turns the earlier expansion layer into a genuine
lemma-development release. The aim is not to replace the existing structural
chapters, but to make the repeated internal moves of the proof architecture
explicit enough that later sections may cite them as stable intermediate
devices. Three burdens are addressed here: the closure spine is split into
reusable intermediate steps, the passage from probes to admissible sectors is
cast as a sequence of named reductions, and the projection discipline is tied
to a small but durable family of proof lemmas.


\begin{lemma}[Triolic admissibility precedes compensator balancing]
\label{lem:triolic-before-hc}
A configuration cannot enter the Heisenberg-compensator balancing regime unless
its triolic decomposition is already admissible with respect to the standing
orthogonality and frame-shift constraints.
\end{lemma}

\begin{proof}
The compensator does not create a triolic decomposition; it only identifies and
balances compatible representatives across already existing sector data. If the
triolic organization were itself degenerate or incompatible with the orthogonal
third direction, then there would be no well-defined data for HC to balance.
Hence triolic admissibility is logically prior to the compensator step.
\end{proof}

\begin{lemma}[Closure filtering follows compensator balancing]
\label{lem:closure-after-hc}
Closure filtering is meaningful only after HC has removed representational drift
between coupled sector descriptions.
\end{lemma}

\begin{proof}
Closure tests are carried out on candidate configurations that are meant to
represent one common admissible state. Without compensator balancing, distinct
representatives may still differ by unresolved phase or sector drift, so a
closure residue would not yet measure physical inadmissibility but only
misalignment of representation. Therefore HC balancing is the prerequisite that
turns closure filtering into a structurally meaningful test.
\end{proof}

\begin{proposition}[Closure ladder]
\label{prop:closure-ladder}
The passage from local structural data to observer-frame physics factors through
the closure ladder described above, and none of its intermediate steps may be
skipped without breaking the proof discipline
\[
\text{closure before projection before interpretation}.
\]
\end{proposition}

\begin{proof}
Lemma~\ref{lem:triolic-before-hc} places admissible triolic organization before
HC balancing. Lemma~\ref{lem:closure-after-hc} places closure filtering after
HC and before any observer readout. The observer frame, however, names only the
output of the stabilized configuration, so projection necessarily comes after
closure filtering. Interpretation is last because it depends on projected
sector naming rather than on raw structural data. Thus the stated factorization
is forced by the standing proof order.
\end{proof}

\subsection*{2. Probe reduction chain and admissible channels}
The probe sector now receives a reusable reduction chain. Instead of speaking
only about a general search space, v1.4.0 fixes the following sequence:
\[
\begin{aligned}
&\text{probe pair}
\Longrightarrow \text{orthogonal search geometry}\\[0.15em]
&\Longrightarrow \text{candidate channel family}
\Longrightarrow \text{tool-guided compression}\\[0.15em]
&\Longrightarrow \text{closure-compatible channel}.
\end{aligned}
\]
This makes explicit where frequency detection, routing compression, Markov
transition guidance, and Newton-type convergence tools actually enter.

\begin{lemma}[Orthogonal probe geometry restricts channel growth]
\label{lem:probe-geometry-restricts}
When two probes are coupled together with the standing orthogonal third group,
the set of admissible candidate channels is a strict restriction of the naive
pairwise interaction space.
\end{lemma}

\begin{proof}
The naive pairwise space allows all pair couplings consistent with the two
incoming probe descriptions. The orthogonal third group imposes an additional
compatibility burden that forbids precisely those channels that would resolve
the pairwise mismatch by leaving the triolic discipline. Hence the admissible
candidate family is smaller than the unrestricted pairwise space.
\end{proof}

\begin{lemma}[Tool-guided compression preserves admissibility order]
\label{lem:tool-compression-order}
Kasisky-, hash-, rainbow-, Markov-, and Newton-type tools are admissible only
when used to compress or navigate the candidate channel family without changing
the structural order of admissibility testing.
\end{lemma}

\begin{proof}
These tools may accelerate pattern detection, routing, transition guidance, or
convergence, but none of them decides admissibility independently of closure.
Their acceptable role is therefore subordinate: they reduce the burden of
search inside the candidate family while the final decision remains the closure
test. Consequently, they preserve admissibility only if they do not replace the
closure filter by a purely computational surrogate.
\end{proof}

\begin{proposition}[Probe reduction principle]
\label{prop:probe-reduction-principle}
A probe channel is physically readable only if it is obtained as the end point
of the probe reduction chain and survives closure filtering after tool-guided
compression.
\end{proposition}

\begin{proof}
By Lemma~\ref{lem:probe-geometry-restricts}, the orthogonal search geometry
already restricts the candidate family. By Lemma~\ref{lem:tool-compression-order},
the technical tools may only compress or navigate that family, not redefine its
criterion of admissibility. Therefore the only physically readable channels are
those that remain after the final closure-compatible selection step.
\end{proof}

\subsection*{3. Projection robustness and controlled interpretation}
The next recurrent burden is the distinction between structural forcing,
projective readout, and phenomenological naming. Version v1.4.0 fixes
this burden by introducing small reusable robustness statements.

\begin{lemma}[Projection does not create structure]
\label{lem:projection-no-creation}
A projection map may translate an already closure-stable structural state into
an observer-frame description, but it cannot create a new admissible structure
that was absent before projection.
\end{lemma}

\begin{proof}
Projection acts on the representative of a state after the admissibility burden
has already been settled. If projection could create admissibility, then the
observer frame would retroactively decide structural existence, contradicting
the standing principle that measurement is phase-shifted relative to the
underlying structure. Thus projection is readout, not creation.
\end{proof}

\begin{corollary}[Interpretation discipline]
\label{cor:interpretation-discipline}
Matter, gauge, mass, Higgs, and geometric language are admissible only as
controlled names for projected closure-stable regimes and not as independent
axiomatic inputs to the structural proof.
\end{corollary}

\begin{proof}
By Lemma~\ref{lem:projection-no-creation}, projection only reads out a state
that has already survived the structural admissibility filter. Therefore the
phenomenological names attached afterwards inherit their legitimacy from the
closure-stable regime and cannot function as prior axioms of the construction.
\end{proof}

\subsection*{4. Technical use inside the monolith}
The practical purpose of these new lemmas is modest but important. They create
stable insertion points for later derivation chains. Subsequent versions may
therefore reference:
\begin{itemize}[leftmargin=2em]
  \item the \emph{closure ladder} when passing from local triolic data to
  observer-frame sectors;
  \item the \emph{probe reduction principle} when compressing search spaces; and
  \item the \emph{projection robustness statements} when interpreting effective
  field sectors.
\end{itemize}
This means that the monolith can deepen individual proofs without reopening the
global explanatory burden every time a local derivation is added.

\subsection*{5. Version status of v1.4.0}
Relative to the previous release, v1.4.0 contributes a stricter internal
proof scaffold. The master-equation ladder and the micro-to-macro chain are
retained, but they are no longer left as merely global narratives. They are now
supported by named intermediate lemmas that later sections may invoke directly.
That is the specific function of v1.4.0: not to conclude the full proof,
but to make the next proof layer technically attachable in a disciplined way.
